\SourceLeaf{251}%
corresponding to the triangle is to the entire surface of the sphere as \(\pm(A+B+C-\pi)\) is to \(4\pi\). This theorem, which, if we are not mistaken, seems to rank among the most elegant in the theory of curved surfaces, can also be stated as follows:

\textit{The amount by which the sum of the angles of a triangle formed by shortest lines on a concave-concave curved surface exceeds \(180^0\), or by which the sum for a triangle formed by shortest lines on a concave-convex curved surface falls short of \(180^0\), is measured by the area of the part of the sphere’s surface corresponding to that triangle through the directions of the normals, if the entire surface area is taken to equal 720 degrees.}

More generally, in any polygon with \(n\) sides, each formed by a shortest line, the excess of the angle sum over \(2n-4\) right angles, or its deficit from \(2n-4\) right angles (according to the nature of the curvature of the surface), equals the area of the corresponding polygon on the sphere’s surface, provided the entire surface of the sphere is taken to equal 720 degrees. This follows directly from the preceding theorem by dividing the polygon into triangles.
\Article{21}
Let us restore to the symbols \(p,q,E,F,G,\omega\) the general meanings they had above, and suppose that the nature of the curved surface is also determined in a similar way by two other variables \(p',q'\), with the general line element expressed by
\[
\BareRoot{(E'\mathrm{d}p'^2+2F'\mathrm{d}p'\mathord{.}\mathrm{d}q'+G'\mathrm{d}q'^2)}
\]
Thus definite values of \(p,q\) specifying any point of the surface correspond to definite values of \(p',q'\). These will therefore be functions of \(p,q\), and we suppose that differentiating them gives
\[
\begin{aligned}
\mathrm{d}p'&=\alpha\mathrm{d}p+\beta\mathrm{d}q\\
\mathrm{d}q'&=\gamma\mathrm{d}p+\delta\mathrm{d}q
\end{aligned}
\]
We now propose to investigate the geometric meaning of these coefficients \(\alpha,\beta,\gamma,\delta\).

\textit{Four} systems of lines can now be conceived on the curved surface, in which \(q,p,q',p'\) are constant, respectively. Suppose we draw through a fixed point, corresponding to the values \(p,q,p',q'\) of the variables, the four lines belonging to the respective systems. The elements of these lines corresponding to the positive increments
